\section{Discussion and limitations}
\label{sec:discussion}

The experiments show that model size is not a reliable predictor of the reported solver time. Some of
the largest instances require very little time inside \texttt{optimize!}, whereas several smaller
instances use most of the limit. The important factor is whether the relaxation produces a lower bound
that is close to a good integer routing. When the two values meet, optimality can be certified without
exploring a large branch-and-bound tree. When they remain far apart, a long search may still leave a
nonzero gap.

The reconfiguration budget makes this issue more visible in the two-period experiment. Although the
implementation combines valid single-period lower bounds with a joint relaxation bound, the resulting
lower bound can still be weak for difficult instances. This is especially relevant for setA-16 and
setA-19, whose mean gaps after one hour are 66.82\% and 55.86\%. Their best feasible MLU values over
the ten runs are 0.1025 and 0.0796, respectively, both noticeably lower than their reported means. The
variation between seeds suggests that additional path generation can still improve the incumbent,
while the large gaps indicate that the current bounds are insufficient to assess how close those
incumbents are to the true optimum.

The averages also hide differences between individual runs. SetA-03 is a clear example: eight runs
prove the same two-period optimum, but two runs reach the time limit with a slightly worse solution.
For this reason, the table marks the instance as not certified rather than treating its small mean gap
as proof of optimality. Similarly, a zero branch-and-bound node count does not mean that no
optimization was performed; the relaxation and column-generation stages may already have closed the
gap at the root. Node counts should also be compared with care because they are summed over solver
calls and come from two different formulations.

The timing results have a further limitation. They follow the course definition and measure only
calls to \texttt{optimize!}; they are not end-to-end running times. This distinction matters most for
column generation, where shortest-path preprocessing, pricing, and construction of new columns take
place outside the solver calls. Since total wall-clock time was not stored in the result files, it
cannot be reconstructed reliably afterwards. The present tables can therefore compare solver effort,
but not the complete runtime of the two solution approaches.

Three limitations define the scope of the conclusions. First, the model minimizes only the maximum
link utilization in~\eqref{eq:mlu}. The full challenge objective is lexicographic: after minimizing the
largest load, it also minimizes the second largest load, then the third largest, and so
on~\citep{orange2026}. Two solutions with the same MLU may therefore have different values under the
official objective. Second, the coefficients $r^t_{ija}$ follow the formula given in the course
note~\citep{courser2026}. This formula assigns equal weight to complete shortest paths and may differ
from router-local ECMP splitting on an asymmetric shortest-path graph. Third, the path-based integer
model contains only the columns generated or added during the run. Its incumbents are feasible for the
original model, and the reported lower bounds remain valid, but a non-certified row is not proof that
the best possible routing has been found.

Future improvements could focus on strengthening the joint relaxation, improving the pricing strategy,
and incorporating branch-and-price so that new paths remain available throughout the integer search.
After the scalar MLU is certified, the same framework could also be extended to optimize the remaining
levels of the lexicographic objective.
